\section{Discussion}
\label{sec:discussion}

\subsection{$w_\text{vra}$ Sensitivity}

The alignment weight $w_\text{vra}$ is the primary tuning parameter of
VRASection.
Its effect on the algorithm's behaviour can be characterised across
three qualitative regimes:

\paragraph{Low weight ($w_\text{vra} \leq 0.2$): compactness dominates.}
At $w_\text{vra} = 0.2$, the maximum alignment discount is 20\,\% of the
normalised edge-cut.
In practice, this means VRASection only deviates from GeoSection's
ratio selection when a highly concentrated split ($A \approx 0.8$--$1.0$)
is nearly as compact as the GeoSection winner.
States with moderate geographic sorting (NC, GA, TX) will show little
change from GeoSection, because their minority populations, while
concentrated in cities, do not produce alignment scores high enough to
overcome a 4:1 compactness preference.
This regime is recommended when legal risk is high and the state has no
clear Section 2 obligation.

\paragraph{Default weight ($w_\text{vra} = 0.4$): balanced.}
The proposed default gives compactness 60\,\% weight and alignment 40\,\%.
For Alabama, this weight is sufficient to shift the ratio from $1:6$ to $2:5$
(Section~\ref{sec:evaluation}).
The Shaw predominance condition is satisfied with a 20\,pp margin.
This regime is recommended as the standard setting for states where
Section 2 litigation has occurred or is anticipated.

\paragraph{High weight ($w_\text{vra} \geq 0.6$): alignment dominates.}
At $w_\text{vra} = 0.6$, the alignment bonus can overcome up to 60\,\% of
the compactness cost, meaning a ratio with $A = 1$ would be preferred
even if its normalised edge-cut is 2.5$\times$ the GeoSection winner's.
In states with moderate geographic concentration, this can produce ratios
that sacrifice meaningful compactness for marginal alignment gains.
The Shaw predominance condition is not satisfied at $w_\text{vra} = 0.6$
(see \S\ref{sec:shaw-line} below).
This regime should not be used in standard redistricting; it may be useful
for research to identify the maximum alignment achievable under geographic
constraints.

\paragraph{Summary table.}
\begin{table}[h]
\centering
\caption{$w_\text{vra}$ regimes and their properties.}
\label{tab:wvra}
\begin{tabular}{ccccl}
\toprule
$w_\text{vra}$ & \textbf{Compactness weight} & \textbf{Alignment weight}
  & \textbf{Shaw compliant?} & \textbf{Recommended use} \\
\midrule
0.0  & 100\,\% & 0\,\%   & Yes (=GeoSection)  & Pure compactness \\
0.2  & 80\,\%  & 20\,\%  & Yes                & Low-obligation states \\
0.4  & 60\,\%  & 40\,\%  & Yes                & \textbf{Default (Section 2 states)} \\
0.5  & 50\,\%  & 50\,\%  & Borderline         & Not recommended \\
0.6  & 40\,\%  & 60\,\%  & No                 & Research only \\
1.0  & 0\,\%   & 100\,\% & No                 & Research only \\
\bottomrule
\end{tabular}
\end{table}

\subsection{The Shaw v. Reno Line}
\label{sec:shaw-line}

The racial-predominance test from \emph{Shaw v.\ Reno} (1993) and
\emph{Miller v.\ Johnson} (1995) asks whether race was the ``predominant
factor'' in district design.
VRASection's objective function makes this question answerable with
a single inequality.

For any ratio $i$ with $\mathrm{EC}_\text{norm}(i) \gg 1$, the selection
score decomposes as:
\[
  \underbrace{(1 - w_\text{vra}) \cdot \mathrm{EC}_\text{norm}(i)}_\text{compactness component}
  \;\;+\;\;
  \underbrace{(-w_\text{vra} \cdot A \cdot \mathrm{EC}_\text{norm}(i))}_\text{alignment component}
\]

The compactness component exceeds the absolute alignment component when:
\[
  (1 - w_\text{vra}) > w_\text{vra} \cdot A \quad\Longleftrightarrow\quad
  w_\text{vra} < \frac{1}{1 + A}
\]

For $A = 1$ (perfect concentration), this requires $w_\text{vra} < 0.5$.
For $A = 0.5$ (moderate concentration), the threshold rises to $w_\text{vra} < 0.67$.
For typical empirical alignment scores ($A \approx 0.3$--$0.5$), the
proposed $w_\text{vra} = 0.40$ satisfies the predominance condition with
a comfortable margin.

This analysis provides a quantitative basis for expert testimony:
``The algorithm assigns geographic compactness a minimum weight of 60\,\%
and alignment a maximum weight of 40\,\%.
For the Alabama 2:5 split, where alignment was 0.42, geographic compactness
received 60\,\% weight and alignment 40\,\%, making geographic compactness
the predominant criterion by a factor of 1.5.''

\subsection{Relationship to B.3: MetisVra Failures}

VRASection and MetisVra address the same goal (VRA-compliant maps) via
fundamentally different mechanisms.
The comparison is instructive.

\paragraph{MetisVra: explicit racial concentration.}
MetisVra adds minority VAP fraction as a second METIS constraint.
The constraint requires each half of the bisection to contain approximately
the target minority fraction.
This approach fails in three of five states because:
(1) METIS's constraint solver cannot handle the 60--800$\times$ scale
mismatch between population balance and minority fraction;
(2) even with extreme tolerance relaxation, the constraint provides
insufficient guidance to overcome geographic dispersion;
(3) the approach likely constitutes racial predominance under \emph{Miller}
in states where achieving the minority target requires sacrificing compactness.

\paragraph{VRASection: geographic signal extraction.}
VRASection uses minority VAP as input to a geographic score, not as a
constraint.
It fails gracefully: when minority population is geographically dispersed,
$A \approx 0$ for all ratios, and the algorithm reduces to GeoSection.
It does not attempt to force a minority-concentrated partition where none
exists in the geography.
This is the correct behaviour: if no geographic boundary separates
minority-concentrated from minority-dispersed areas, there is no
Gingles Prong 1 community to protect, and a standard compact-district
algorithm is the legally appropriate choice.

\paragraph{Zero user-specified minority thresholds.}
MetisVra requires a user-specified minority target fraction (e.g., $>50\,\%$)
and tolerance (ubvec).
These parameters are legally fraught: a practitioner who specifies
``50\,\% minority target'' is making an explicit racial decision.
VRASection requires no such parameters.
The $w_\text{vra}$ weight is a methodological preference
(compactness vs.\ alignment balance), not a racial target.
Two practitioners using VRASection with $w_\text{vra} = 0.40$ will
produce the same map, without either having specified a minority fraction.

\subsection{Limitation: Alignment Does Not Guarantee $> 50\,\%$ Minority VAP}

VRASection's alignment score measures geographic concentration at the
\emph{bisection level}, not at the \emph{district level}.
A high-alignment first-level split concentrates minority population on
one side, improving the conditions for MM district formation in the
downstream recursion.
But it does not guarantee that any specific district will exceed 50\,\%
minority VAP.

This limitation is inherent to any first-level-only VRA intervention.
The full recursion is conducted by standard GeoSection, which optimises
compactness without reference to minority population.
In the southern sub-region of Alabama's $2:5$ split, the 35--40\,\%
Black VAP must be further concentrated by the downstream bisection tree.
Whether this produces one or two Black-majority districts depends on the
detailed geography of the Black Belt sub-region and the specific ratio
selections in subsequent bisections.

A post-hoc audit is recommended: after running VRASection, compute
minority VAP for every final district and report the count of MM districts.
If the count falls below the Gingles target, the analyst should consider
increasing $w_\text{vra}$ or applying VRASection at lower recursion levels
(a ``VRA-aware full recursion'' mode discussed in the plan.md open questions
but not recommended for standard use due to increased racial predominance risk).

\subsection{Callais v. Landry and the Evidence Standard}

The redist codebase's Callais evidence layer (landed 2026-04-30) provides
WLS+HC3 regression infrastructure for validating racial bloc voting
under the \emph{Callais v.\ Landry} (2025) standard.
VRASection is designed to be used in conjunction with this evidence layer:
the algorithm produces a map candidate; the bloc-voting analysis validates
whether each candidate minority opportunity district satisfies Gingles
Prongs 2 and 3 in the jurisdiction.

This two-stage workflow --- VRASection for map generation, bloc-voting
analysis for evidence production --- separates the algorithmic decision
(geographic concentration) from the statistical evidentiary question
(racially polarised voting).
The separation is legally significant: the map drawer can state that
district boundaries were set by geographic compactness (VRASection),
and the VRA claim is supported by independent statistical evidence
(Callais analysis), with no circular dependency between the two.
